\documentclass[12pt]{article}
\usepackage[utf8]{inputenc}
\usepackage[T1]{fontenc}
\usepackage{hyperref}
\title{Wangkangurru Language Resources for K-12 Curriculum Development}
\author{Compiled by Claude Sonnet 5.5 (claude-sonnet-5-5) for Andr\'as Kornai}
\date{}
\begin{document}
\maketitle

\section{Introduction}
Wangkangurru (ISO 639-3 wgg; Glottocode wang1290) is a Pama-Nyungan language of Australia, associated with the Birdsville region. It is not spoken on a daily basis. The Mobile Language Team describes it as critically endangered, and some sources list it as extinct. The Wangkangurru/Yarluyandi community keeps cultural and intellectual ownership of the language and asks to be consulted before it is used outside the community. Any school curriculum should be built in consultation with that community. The searches found a small set of resources, listed below.

\section{Primary Educational Resources}
The Mobile Language Team's Online Language Learning Portal has lessons developed with the Crombie family, who are speakers of Eastern Wangkangurru. The lesson page is at \url{https://portal.mobilelanguageteam.com.au/?p=4251}. The Mobile Language Team's Wangkangurru page is at \url{https://mobilelanguageteam.com.au/languages/wangkangurru/}.

\section{Linguistic and Reference Materials}
\emph{Wangkangurru wangka-purru} (2022) is a community dictionary compiled by Eleanor McCall with the Crombie family and published by the Mobile Language Team. It has an introduction, a pronunciation guide, phrases, stories, and English--Wangkangurru and Wangkangurru--English word lists. It is available as a PDF at \url{https://portal.mobilelanguageteam.com.au/wp-content/uploads/2022/05/Wangkangurru-dictionary-2022.pdf}.

Luise Hercus, \emph{A Grammar of the Arabana-Wangkangurru Language} (Pacific Linguistics, 1994), is reported to be available for download from the Mobile Language Team site. No direct URL for the grammar appeared in the search results. A Hercus draft dictionary is mentioned in the 2022 dictionary's introduction. It was not found as a published or downloadable resource.

A general overview is at \url{https://en.wikipedia.org/wiki/Wangkangurru_dialect}.

\section{Religious and Community Resources}
No scripture or religious materials were found. The community organisation behind the language work is the Mobile Language Team, which is based in Adelaide: \url{https://mobilelanguageteam.com.au/?p=223}.

\section{Digital and Audio Resources}
The online lesson portal listed above is the only digital resource found. No audio or video collections were found in the search results.

\section{Curriculum Development Context}
The language is not known to be taught in schools. The 2022 dictionary, the portal lessons and the Hercus grammar are the basis for a heritage and revitalization curriculum. Work should be done with the Wangkangurru/Yarluyandi community and the Mobile Language Team. The Creative Commons terms on the dictionary should be checked, along with the community's request to be consulted. Little else was found, so the curriculum would start from a small base.
\end{document}